\section*{Proposed Solution} \label{sec:solution}
\noindent
MECC addresses the four unmet requirements through a two-part invention.

\subsection*{Part 1: Feistel-Permuted Grid Hash ECC (Primary Invention --- 80\%)}

The core of MECC is a bit-flip detection and correction pipeline comprising three stages.

\textbf{Stage 1 --- Feistel Permutation.}
Given a flat bit array of length $L$ and a secret key $k$, a cycle-walking Feistel cipher maps each source index $i \in [0, L)$ to a unique grid coordinate $(r, c)$ in an $N \times N$ grid, where $N = \lceil\sqrt{L}\rceil$.
The permutation uses SHA-256 as the round function.
The key property is that adjacent source indices (which would form a burst in the original sequence) are mapped to geometrically distant grid positions, distributing burst errors uniformly and randomly across the grid.
This equalization means the solver sees effectively random errors regardless of the original error pattern.

\textbf{Stage 2 --- CRC Hash Graph.}
Non-overlapping groups of consecutive rows and groups of consecutive columns of the permuted grid are each hashed using a CRC checksum (CRC-8, CRC-16, or CRC-32, selectable as a design parameter).
Row-group hashes and column-group hashes together cover every cell in the grid: cell $(r, c)$ belongs to exactly one row-group hash and one column-group hash.
Pairs of row and column hash nodes with overlapping source-index coverage are connected by undirected edges, weighted by the number of shared source bits, forming a bipartite hash graph.
This graph encodes which hash nodes share source bits and can therefore help disambiguate each other during correction.

\textbf{Stage 3 --- Coverage-Ordered Solver.}
During correction, hash nodes are processed in order of increasing coverage (smallest number of source bits first).
For each mismatched node, the solver enumerates all combinations of 0, 1, 2, \ldots flip counts within the node's covered source indices, evaluating each candidate by re-hashing and counting total matched nodes across the whole grid.
The candidate that maximizes total matched nodes is accepted.
Because small groups are processed first, the search space is narrow at each step, and corrections propagate forward to adjacent nodes through the graph edges.

\textbf{Configurable overhead and scaling.}
The total overhead is $(G_r + G_c) \cdot B / L$ bits, where $G_r$ and $G_c$ are the number of row and column groups and $B$ is the CRC width in bits.
With groups of fixed size, $G_r \approx G_c \approx \sqrt{L}$, so overhead scales as $O(B\sqrt{L}/L) = O(1/\sqrt{L})$.
As block size doubles, overhead fraction decreases by $\approx 30\%$.
The CRC width and group size are independently tunable, allowing practitioners to dial in the desired tradeoff between correction capacity and storage cost.

\subsection*{Part 2: Embedded Metadata Storage (Complementary Invention --- 20\%)}

When the protected data consists of structured numeric values with representational degrees of freedom --- quantized neural network weights being the canonical example --- the CRC hash metadata need not be stored separately at all.

The embedded variant co-optimizes the data values and the CRC bits jointly via a mixed-integer linear program (MILP):

\begin{center}
\textit{Minimize} $\displaystyle\sum_{i=1}^{W} |v_i - v_i^{\text{eff}}|$ \quad \textit{subject to:} CRC constraints are satisfied by the bit representation of $v_i^{\text{eff}}$.
\end{center}

The MILP finds a data value assignment $\{v_i^{\text{eff}}\}$ that satisfies all CRC parity constraints while minimizing the total distortion from the original values $\{v_i\}$.
The result is that the CRC hash bits are ``hidden'' within small perturbations to the payload values, and the full correction capability of the hash graph is preserved.
Effective storage overhead drops from $O(1/\sqrt{L})$ to near zero.

A \textbf{hybrid scheme} partitions the CRC groups: some groups use embedded storage (zero overhead, small distortion), while groups covering the most sensitive data use explicit storage (small overhead, zero distortion).
The fraction assigned to each mode is a configurable hyperparameter.
This hybrid approach extends the range of achievable (overhead, distortion) operating points well beyond what either pure strategy can achieve.

\textbf{Combined benefit.}
The two-part invention enables a range of deployments: high-capacity error correction with $O(1/\sqrt{L})$ overhead for general-purpose storage, and near-zero overhead for distortion-tolerant applications such as edge neural network inference, where the payload naturally accommodates small value perturbations.
